% --- Archon physics formalization source begin ---
% archon:physics
% archon:covers IPhO2026Problems/problem_ipho_2026_t3_b1.lean
% archon:source-report reports/ipho_2026/problem_ipho_2026_t3_b1.source.json
% archon:problem-id ipho_2026_t3
% archon:part-id T3-B1
% archon:previous-part ipho_2026_t3_a3
% archon:previous-part-policy derive_inline_from_problem_only_material

\chapter{Ipho Physics problem ipho\_2026\_t3\_b1}
\label{ch:IPhO2026Problems_problem_ipho_2026_t3_b1}

\paragraph{Problem source.}
\#\# Physical scenario

Continue with the paramagnetic torus.  Its equation of state is T*M*V = n*K*H,
its heat capacity at constant M is C\_M = n*lambda/T\textasciicircum{}2, and dU = C\_M*dT.
The volume is fixed and the magnetic work on the material is
dW = mu\_0*V*H*dM.  Work and heat entering the torus are positive.

\#\# Current subquestion T3-B1

At fixed temperature T, H changes from H\_i to H\_f. Find the heat Q transferred into the torus.

\paragraph{Shared problem context.}
Continue with the paramagnetic torus.  Its equation of state is T*M*V = n*K*H,
its heat capacity at constant M is C\_M = n*lambda/T\textasciicircum{}2, and dU = C\_M*dT.
The volume is fixed and the magnetic work on the material is
dW = mu\_0*V*H*dM.  Work and heat entering the torus are positive.

\paragraph{Current subquestion.}
At fixed temperature T, H changes from H\_i to H\_f. Find the heat Q transferred into the torus.

\paragraph{Figure/image paths.}
\begin{itemize}
\item ipho\_2026\_source/image/T3\_page-3.png
\item ipho\_2026\_source/image/T3\_page-2.png
\end{itemize}

\paragraph{Reusable previous-part conclusions.}
\begin{itemize}
\item Source T3-A3. Question: Subtract the vacuum-core contribution and write the work dW done on the paramagnetic material. Policy: derive\_inline\_from\_problem\_only\_material
\end{itemize}

\paragraph{Formalization target.}
create a compiling Lean file with sorry bodies at `IPhO2026Problems/problem\_ipho\_2026\_t3\_b1.lean`.
The Lean declarations must preserve the physical quantities, dimensions or dimensional roles, figure labels, governing-law hypotheses, and final relation expressed by this problem.
Use Mathlib/Physlib names found through LeanExplore where available. If a domain API is missing, introduce faithful local abstractions rather than scalar placeholder aliases.
This is an answer-blind solve-phase task. Derive a candidate result only from the problem-side evidence above; do not assume a target value, select a tolerance around a desired value, or encode a desired result into a candidate domain.
For numerical questions, define the raw end-to-end quantity and apply an explicit source-derived rounding rule. For identification questions, characterize or prove uniqueness before recording the derived candidate.

\begin{theorem}[Ipho Physics formalization target]
\label{thm:physics:ipho_2026_t3_b1:target}
\leanok
The assigned autoformalize agent should translate this IPhO physics problem into Lean declarations in the covered file, with theorem and lemma proof bodies written as `by sorry`.
\end{theorem}
\begin{proof}
\leanok
Prover stage complete (iteration 001): all four proof holes
(`isothermal_heat`, `isothermal_heat_neg\_iff`,
`heat\_flows\_out\_when\_magnetized`, `heat\_flows\_in\_when\_demagnetized`)
are closed with kernel-checked proofs; the file compiles warning-free and
every theorem depends only on the standard axioms
(\texttt{propext}, \texttt{Classical.choice}, \texttt{Quot.sound}).
\end{proof}
% --- Archon physics formalization source end ---
